% tech-cv.tex
% Maria Szegedy 2024
% below is original header from back in like 2012 when i downloaded the template
% this is distantly based on. minus the warning to use res.cls bc i rewrote this
% with \documentclass{article} like a million years ago

%%%%%%%%%%%%%%%%%%%%%%%%%%%%%%%%%%%%%%%%%
% Medium Length Graduate Curriculum Vitae
% LaTeX Template
% Version 1.1 (9/12/12)
%
% This template has been downloaded from:
% http://www.LaTeXTemplates.com
%
% Original author:
% Rensselaer Polytechnic Institute (http://www.rpi.edu/dept/arc/training/latex/resumes/)
%
%%%%%%%%%%%%%%%%%%%%%%%%%%%%%%%%%%%%%%%%%

%----------------------------------------------------------------------------------------
%   PACKAGES AND OTHER DOCUMENT CONFIGURATIONS
%----------------------------------------------------------------------------------------

\documentclass{article}

\usepackage[top=2cm,bottom=2cm,left=2cm,right=2cm]{geometry}
\usepackage{fontspec}
\usepackage{amsmath}
\usepackage{enumitem}
\usepackage{titlesec}
\usepackage{tabularx}
\usepackage{multirow}
\usepackage[dvipdfmx]{graphicx}

\pagestyle{empty}
\setlength\parindent{0pt}
\linespread{1.15}
\setitemize{noitemsep}
\titlespacing*{\section}
  {0pt}   % left margin
  {1em} % top margin
  {0.7em}   % bottom margin

\setmainfont[
  Path           = ./fonts/ ,
  UprightFont    = * ,
  SmallCapsFont  = *Caps ,
  AutoFakeSlant  = 0.15 ,
  AutoFakeBold   = 1.5 ,
  Extension      = .otf ,
  Ligatures      = TeX ,
  Scale          = 1.2
]{ScalaSans}

\newcommand{\jobentry}[4]{%
  {\it #1} \hfill #3 \\
  #2

  #4
}

\begin{document}

%----------------------------------------------------------------------------------------
%   NAME AND ADDRESS SECTION
%----------------------------------------------------------------------------------------

\begin{tabularx}{\textwidth}{>{\centering\arraybackslash}X@{}r}
  {\large\bf Maria Szegedy} &
  \multirow{6}{*}{
    \begin{tabular}{c}
      See also my online CV: \\
      \includegraphics[width=3.6cm]{qr.eps} \\
      {\fontsize{12}{12} \texttt{https://mszegedy.github.io/cv}}
    \end{tabular}
  } \\[0.6em]\cline{1-1}\vspace{0em}
  {60 Franklin Drive} \\
  {Hillsborough, New Jersey 08844} \\
  {+1 (908) 642-4336} \\
  {\fontsize{12}{12}%
    $\underbrace{\texttt{mszegedy}}%
    _{\textstyle\text{\makebox[0pt]{\fontsize{10}{10}\selectfont Github, Gitlab}}}%
    \texttt{@tutanota.com}$
  }
\end{tabularx}

%----------------------------------------------------------------------------------------

%----------------------------------------------------------------------------------------
%   EDUCATION SECTION
%----------------------------------------------------------------------------------------

\section*{EDUCATION \& CERTIFICATIONS}

{\it AWS Certified Solutions Architect---Associate} \\
\hspace*{3em} Issued 2024 May 10 \\
\hspace*{3em} Expires 2027 May 10 \\
{\fontsize{9}{10}\selectfont
  Available at
  \texttt{https://mszegedy.github.io/cv/assets/pdfs/aws-solutions-architect-assoc-cert.pdf}.}

\vspace{0.8em}

{\it Bachelor of Science,} in Biomedical Engineering (Tissue Engineering and
Molecular Bioengineering option) \\
\hspace*{3em} Additional major in Molecular Biology and Biochemistry \\
\hspace*{3em} Additional major in Physics (professional option) \\
Rutgers University, New Brunswick, NJ, August 2021

\section*{PUBLICATIONS}

Leman,~J.K., Weitzner,~B.D., Lewis,~S.M. \textit{et al}. Macromolecular modeling
and design in Rosetta: new methods and frameworks. \textit{Nat Methods}
\textbf{17}, 665–680 (2020). doi: 10.1038/s41592-020-0848-2.

%----------------------------------------------------------------------------------------
%   PROFESSIONAL EXPERIENCE SECTION
%----------------------------------------------------------------------------------------

\section*{EXPERIENCE}

\jobentry{Research Assistant}
  {Rutgers University, Department of Quantitative Biomedicine, Piscataway, NJ}
  {2017--2020}{
    \begin{itemize}
      \item Research assistant to Professor Sagar Khare.
      \item Used PyRosetta to investigate thermostability of certain designed
        proteins.
      \item Wrote API to allow use of molecular dynamics within Rosetta; presented
        it at Rosetta\-Con 2018 and 2019.
    \end{itemize}
  }

\jobentry{Programming Instructor \& Tutor}
  {Rutgers IEEE N2E Coding Club, Rutgers University}
  {Fall 2019--Spring 2020}{
    \begin{itemize}
      \item Gave weekly lessons on MATLAB to first years for N2E Coding Club.
      \item Tutored biomedical engineering graduate students on Mathematica.
    \end{itemize}
  }

\section*{RELEVANT COURSES}
{\it Intro to Computing for Engineers:} \\
  MATLAB. Received A. \\[0.2em]
{\it Computer-based Experimentation in Physics:} \\
  Signal-processing benchwork and computation. Received A. \\[0.2em]
{\it Modern Physics Experimentation:} \\
  Intense, around-the-clock work with complex physics workflows and equipment. Received B+.

\end{document}
